\documentclass[12pt,a4paper]{article}

\usepackage[utf8]{inputenc}
\usepackage[T1]{fontenc}
\usepackage{lmodern}
\usepackage[margin=1in]{geometry}
\usepackage{setspace}
\usepackage[hidelinks]{hyperref}
\usepackage{amsmath,amssymb}
\usepackage{enumitem}
\usepackage{booktabs}
\usepackage{titlesec}
\usepackage[stretch=15]{microtype}
\usepackage{parskip}

% Section formatting
\titleformat{\section}{\large\bfseries}{}{0em}{}
\titleformat{\subsection}{\normalsize\bfseries}{}{0em}{}
\titleformat{\subsubsection}{\normalsize\itshape}{}{0em}{}

\title{Technical Note: Review of the Literature Section\\[6pt]
\large ``Observable-Space Prior Specification for Bayesian Hierarchical
Models''\\[6pt]
\normalsize Current state on branch \texttt{feat/joe-restructure}}
\author{}
\date{February 2026}

\begin{document}
\maketitle
\thispagestyle{empty}

\vspace{1em}
\hrule
\vspace{1em}

\section{Overview}

This note reviews the literature section (Section~2) of the paper following its
recent rewrite. The section was substantially expanded from the previous
version, which had a concise, journal-article style with four subsections
(Camp~1, Camp~2, BHMs in Economics, Positioning). The new version has two
subsections (Camp~1, Camp~2) with considerably more text and a different
rhetorical strategy.

The note is organized as follows: \S1 assesses the structural changes, \S2
reviews Camp~1, \S3 reviews Camp~2, \S4 identifies missing content from the
previous version, \S5 catalogues writing issues, and \S6 flags bibliographic
problems.


\section{1.\ Structural Assessment}

\subsection{What was gained}

The rewrite makes several improvements:

\begin{itemize}[leftmargin=1.5em]
  \item The new introductory paragraphs (lines 215--230) give a clear,
    accessible framing of the two-camp distinction and explicitly state which
    camp the paper contributes to.
  \item Camp~1 is more thorough: the Jeffreys prior gets a proper paragraph
    explaining invariance and improperness, the IG($\varepsilon$,
    $\varepsilon$) pathology is described in more detail (with a footnote
    tying it to Gao et al.), and the Gelman, Simpson \& Betancourt (2017)
    reference is added.
  \item Camp~2 is significantly expanded with a substantial survey of the
    belief elicitation literature (proper scoring rules, risk-neutralization,
    Becker--DeGroot--Marschak, etc.).
\end{itemize}

\subsection{What was lost}

Two subsections from the previous version were removed entirely:

\begin{enumerate}[leftmargin=1.5em]
  \item \textbf{``BHMs in Experimental and Behavioral Economics''} ---
    Previously cited Gao et al.\ (2022), Nilsson et al.\ (2011), Vincent
    (2016), Del Negro \& Schorfheide (2008), and Villani (2009). This
    provided context for BHMs in the paper's target field and connected to
    macroeconometrics precedents for observable-space priors. All five
    citations are now absent from the paper.
  \item \textbf{``Positioning''} --- Previously stated the paper's
    contribution sharply: ``the first systematic derivation of hyperpriors from
    individual-level beliefs about observable quantities, combining Camp~2's
    predictive insight with Camp~1's analytical rigor.'' This had a clear,
    assertive claim of novelty that the new closing paragraphs (lines
    357--367) do not replicate with the same force.
\end{enumerate}

\subsection{Assessment}

The BHMs subsection should be restored in some form. The Del Negro \&
Schorfheide and Villani references were genuinely useful for positioning ---
they show the idea of observable-space priors has appeared in
macroeconometrics, establishing that the contribution is the systematic
derivation for hierarchical settings, not the idea itself. The Nilsson and
Vincent references are needed because the paper claims three applications
(RDU, CPT, discounting), and these are the primary BHM references in those
fields.

The Positioning subsection's assertive claim of novelty should be
restored, either as a standalone paragraph at the end of Camp~2 or as a
separate subsection. The current closing (lines 357--367) is too tentative ---
it describes what the paper does procedurally rather than claiming what is new.


\section{2.\ Camp 1: Detailed Review}

\subsection{Content}

The Camp~1 subsection is solid. It correctly covers:
\begin{itemize}[leftmargin=1.5em]
  \item Jeffreys/reference priors and their properties
  \item The IG($\varepsilon$, $\varepsilon$) pathology (Gelman 2006)
  \item Half-Cauchy priors (Polson \& Scott 2012)
  \item Horseshoe prior (Carvalho et al.\ 2010)
  \item PC priors (Simpson et al.\ 2017)
  \item Gelman, Simpson \& Betancourt (2017) on prior selection when models
    are developed while examining data
\end{itemize}

\subsection{Issues}

\begin{enumerate}[leftmargin=1.5em]
  \item \textbf{``observations grow'' (line 245):} Should be ``observations
    grows'' (number\ldots grows) or rephrase to ``as the number of observations
    increases.''

  \item \textbf{``while trying to still minimize'' (line 267--268):} Awkward
    split infinitive. Rewrite to ``while still trying to minimize'' or
    ``while minimizing.''

  \item \textbf{``domain specific knowledge'' (lines 275, 278, 280):}
    Should be hyphenated: ``domain-specific knowledge.'' Appears three times
    in the section.

  \item \textbf{Loss of the ``gap'' motivation (lines 275--282):} The old
    version's closing paragraph for Camp~1 was: ``These contributions solve
    genuine problems\ldots However, they share a limitation: they do not
    provide a mechanism for encoding substantive domain knowledge when it
    exists. A researcher who knows from decades of experiments that CRRA
    coefficients typically fall between 0.2 and 0.8 gains nothing from a
    prior designed to express ignorance.'' This was rhetorically effective: it
    acknowledged the value of Camp~1 while sharply motivating Camp~2. The new
    closing paragraph is more diplomatic but less pointed. Consider restoring
    the specific CRRA example --- it is concrete and memorable.

  \item \textbf{Previous version said reference priors ``do not naturally
    extend to hierarchical settings'':} This was a useful observation that
    specifically motivated the paper's hierarchical contribution. The new
    version drops it. Consider restoring.
\end{enumerate}


\section{3.\ Camp 2: Detailed Review}

\subsection{Content}

The Camp~2 subsection covers three threads:
\begin{enumerate}[leftmargin=1.5em]
  \item Predictive elicitation (Kadane 1980, Kadane \& Wolfson 1998)
  \item Belief elicitation in economics and psychology (Juster, Winkler,
    Murphy, Savage, Becker--DeGroot--Marschak, Harrison \& collaborators,
    Offerman, Hossain \& Okui, Smith, Schlag, Garthwaite, Manski, O'Hagan,
    Mikkola)
  \item Prior predictive methods (Gabry et al., Bland, Manderson, Segert
    et al.)
\end{enumerate}

\subsection{Major structural concern: the belief elicitation digression}

The belief elicitation material (lines 297--335) --- proper scoring rules,
risk-neutralization via utility payments, BDM auctions --- is a substantial
digression that is not directly relevant to the paper's contribution.

The paper does not elicit beliefs from experimental subjects. It helps
\textit{researchers} encode \textit{their own} prior beliefs about model
parameters into hyperprior distributions. The belief elicitation literature is
about extracting subjective probabilities from lay individuals in
experimental settings. These are related (both involve turning beliefs into
distributions) but they are different problems, and the connection is not
drawn explicitly.

\textbf{Recommendation:} Either (a) cut the belief elicitation survey to 2--3
sentences acknowledging its existence and citing Schlag (2015) or Manski
(2004) as entry points, or (b) explicitly draw the connection: ``The same
cognitive challenges that arise in eliciting beliefs from experimental
subjects --- difficulty reasoning about probabilities, sensitivity to framing
--- also arise when asking researchers to specify hyperprior parameters.
The observable-space approach developed here addresses this by translating
the specification task into domain quantities (dollars, probabilities) rather
than abstract distribution parameters.''

Without either treatment, the 2.5-paragraph survey reads as showing breadth
of knowledge rather than serving the argument.

\subsection{Issues}

\begin{enumerate}[leftmargin=1.5em]
  \item \textbf{``a substantial literature onto itself'' (line 297--298):}
    Should be ``unto itself'' or ``in itself.''

  \item \textbf{``has lead to techniques'' (line 312):} Should be ``has led
    to techniques.''

  \item \textbf{``Drawing connections\ldots are actively being explored''
    (line 346--347):} Subject-verb disagreement. ``Drawing connections'' is a
    gerund phrase (singular). Should be ``Drawing connections\ldots is actively
    being explored'' or rephrase: ``Connections between prior predictive
    distributions and prior distributions are actively being explored.''

  \item \textbf{Missing contrast with Bland (2025) (lines 342--344):} The
    previous version explicitly characterized Bland's approach as ``informal,
    relying on visual inspection of prior predictive distributions and manual
    adjustment, where the framework developed here provides closed-form
    solutions.'' The new version just says Bland ``proposed calibrating priors
    for structural models by examining their implications for certainty
    equivalents'' without any contrast. The contrast is the whole point of
    citing Bland --- it is the closest precursor. Restore it.

  \item \textbf{Missing contrast with Manderson (2026) (lines 348--350):}
    Same issue. The previous version noted Manderson ``requires black-box
    numerical optimization where the framework developed here provides
    closed-form solutions.'' The new version just describes what Manderson
    does. The contrast is needed to establish the contribution.

  \item \textbf{Missing characterization of Segert et al.\ (2019) (lines
    351--355):} The previous version noted that Segert et al.\ ``established
    the mathematical foundation that the present work extends to hierarchical
    settings.'' The new version mentions they focus on non-informative
    predictive distributions but drops the claim about extending their
    foundation. This was a useful positioning statement.

  \item \textbf{Closing paragraphs (lines 357--367):} ``There are many
    threads being pulled'' is colloquial and vague. ``Far less insight about
    what properties should be held by the hyperpriors'' is passive and
    indirect. Compare to the old Positioning subsection's claim of ``the first
    systematic derivation.'' The current closing does not make a clear novelty
    claim. It should.

  \item \textbf{``the literature of risky choice where expert beliefs about
    the prior predictive of risk preferences are well-informed'' (lines
    363--364):} This is imprecise. It is not the ``prior predictive of risk
    preferences'' that is well-informed; it is researchers' beliefs about
    observable quantities (CEs, choice probabilities) that are well-informed
    by decades of experimental evidence. Clarify.

  \item \textbf{``I begin by formally connecting'' (line 364):} This reads as
    a transition to the next section but is placed at the end of the
    literature review. It would fit better as the opening of Section~3, or
    reworded as a statement of contribution rather than a procedural
    announcement.
\end{enumerate}


\section{4.\ Missing Content from Previous Version}

The following content was present in the previous version and should be
restored or consciously replaced:

\begin{enumerate}[leftmargin=1.5em]
  \item \textbf{BHMs in experimental economics subsection:}
    \begin{itemize}
      \item Gao et al.\ (2022) as the specific comparison case (mentioned
        in Camp~1 footnote but not positioned as the comparison benchmark)
      \item Nilsson, Rieskamp \& Wagenmakers (2011) --- hierarchical CPT,
        directly relevant to the planned CPT application
      \item Vincent (2016) --- hierarchical discounting, directly relevant to
        the planned discounting application
    \end{itemize}

  \item \textbf{Observable-space priors in macroeconometrics:}
    \begin{itemize}
      \item Del Negro \& Schorfheide (2008) --- DSGE priors from
        model-implied moments matching pre-sample data
      \item Villani (2009) --- steady-state priors for VARs, beliefs about
        observables rather than coefficients
    \end{itemize}
    These were valuable because they show the idea has precedent in other
    fields, which strengthens the positioning.

  \item \textbf{Explicit novelty claim:} ``The present paper provides the
    first systematic derivation of hyperpriors from individual-level beliefs
    about observable quantities, combining Camp~2's predictive insight with
    Camp~1's analytical rigor.''

  \item \textbf{$\rho$ and $\kappa$ comparison to PC priors:} ``The $\rho$
    and $\kappa$ parameters provide the same interpretability as PC priors'
    threshold parameters while encoding positive beliefs rather than
    penalizing complexity.'' This was a sharp and useful comparison.
\end{enumerate}


\section{5.\ Writing Issues (Catalogue)}

In addition to the issues noted in context above:

\begin{enumerate}[leftmargin=1.5em]
  \item Line 216: ``will result in proper posterior distribution'' --- missing
    article: ``will result in \textit{a} proper posterior distribution.''
  \item Line 238: ``\textbackslash enquote\{let the data speak.\}'' --- the
    period should be outside the closing quote in British convention, or inside
    in American. Be consistent with the rest of the paper. The previous
    version used ``let the data speak'' with double-backtick quotes.
  \item Line 293--294: ``\textbackslash enquote\{structural elicitation\}''
    uses enquote but then ``\textbackslash textit\{predictive elicitation\}''
    uses italics. Pick one convention for defined terms.
  \item Lines 340--341: ``they suggest that prior predictive checking be used
    iteratively to check the reasonableness'' --- ``check\ldots check'' is
    repetitive. Rephrase: ``\ldots be used iteratively to assess the
    reasonableness.''
  \item Line 348: ``\textcite\{Manderson2026\}'' --- Verify this date. The
    paper was a 2023 preprint. If it has been published in 2026, confirm the
    venue and update the bib entry. If it is still a preprint, cite it as
    such.
  \item Line 351: ``\textcite\{Segert2019\}'' --- The previous version cited
    this as ``SegertDavisStober2019.'' Verify the bib key is correct and that
    the entry exists.
\end{enumerate}


\section{6.\ Bibliographic Issues}

\begin{enumerate}[leftmargin=1.5em]
  \item \textbf{Bland2025 is undefined.} It is cited at lines 181 and 342 but
    has no entry in \texttt{doc/library.bib}. The previous version cited it
    as Bland2023. Either the paper has been published (add the bib entry) or
    the date was changed speculatively (revert to the correct date and add the
    entry).

  \item \textbf{Missing bib entries from removed content.} The following
    citations from the previous version are no longer in the paper and their
    bib entries may not exist in \texttt{library.bib}:
    \begin{itemize}
      \item Vincent2016
      \item DelNegroSchorfheide2008
      \item Villani2009
    \end{itemize}
    If the BHMs subsection is restored, these entries will need to be added.

  \item \textbf{Citation key changes.} The previous version used
    ``SegertDavisStober2019''; the new version uses ``Segert2019''. Verify
    which key actually exists in the bib file(s).

  \item \textbf{Manderson date.} Previously cited as Manderson2023, now as
    Manderson2026. Verify publication status.

  \item \textbf{Bland date.} Previously Bland2023, now Bland2025. Verify.

  \item \textbf{Many new citations lack bib entries.} The new Camp~2 content
    cites Juster1966, Winkler1969, Murphy1970, Savage1971, Becker1964,
    Harrison2017b, Harrison2019a, Offerman2009, Hossain2013, Smith1961,
    Schlag2015, Manski2004, Manski2010. These all exist in
    \texttt{library.bib} --- confirmed. But the entries should be
    cross-checked for completeness (volume, pages, DOI).
\end{enumerate}


\section{7.\ Summary of Recommendations}

In priority order:

\begin{enumerate}[leftmargin=1.5em]
  \item \textbf{Fix Bland2025 bib entry} --- the paper won't compile with
    biber until this is resolved.
  \item \textbf{Restore the BHMs subsection} (or fold the key references into
    Camp~2) --- Nilsson, Vincent, Del Negro \& Schorfheide, and Villani are
    needed for positioning.
  \item \textbf{Restore explicit contrasts} with Bland, Manderson, and Segert
    et al.\ --- without these, the novelty claim is implicit rather than
    explicit.
  \item \textbf{Restore or rewrite the Positioning paragraph} --- the paper
    needs a clear, assertive claim of what is new.
  \item \textbf{Cut or connect the belief elicitation survey} --- either
    shorten it substantially or draw the explicit connection to the paper's
    contribution.
  \item \textbf{Fix the grammatical and writing issues} catalogued in \S5.
  \item \textbf{Resolve citation key discrepancies} (Segert, Manderson, Bland
    dates/keys).
\end{enumerate}

\end{document}
